\documentclass[11pt,a4paper]{report}

\usepackage{float} %dzieki temu pakietowi rysunki pojawiaja sie tam gdzie chce po uzyciu opcji H w srodowisku figure
\usepackage[font=small,labelfont=bf]{caption}
\usepackage{natbib}
\usepackage{import}
\usepackage[tight]{subfigure}
\usepackage{tikz}
\usepackage{pgfplots}
\usepackage{polski}
\usepackage[cp1250]{inputenc}
\usepackage{indentfirst}
\usepackage{color}
\usepackage{graphics}
\usepackage{geometry}
% ---- PAGE LAYOUT ----
\geometry{a4paper, top=1.5cm, bottom=1.5cm, left=1.5cm,
right=1.5cm, nohead}
\usepackage{amsmath,amssymb,amsfonts}
\begin{document}
\begin{center}
\large{\textbf{PRACA DOMOWA II}}
\end{center}

\noindent
imi� i nazwisko ........................................\\


\noindent
\textbf{Zadanie 1} Oblicz ca�ki podw�jne przechodz�c do wsp�rz�dnych biegunowych $r,\varphi$:\\
\textbf{(a)} $\iint\limits_{D} (x^2+y^2) dxdy$, gdzie $D$: $x^2+y^2\leq 4$, $y\geq -x$, $x\geq 0$,\\
\textbf{(b)} $\iint\limits_{D} x y^2 dxdy$, gdzie $D$: $1\leq x^2 + y^2 \leq 4$, $y\geq 0$, $y\leq x$,\\
\textbf{(c)} $\iint\limits_{D} \frac{x y dxdy}{x^2+y^2 + 1}$, gdzie $D$: $x^2 + y^2 \leq 1$, $y\geq \sqrt{3}x$, $x\geq 0$,\\
\textbf{(d)} $\iint\limits_{D} xy dxdy$, gdzie $D$: $x^2+y^2\leq 2x$, $y\geq 0$\\

\vspace{0.25cm}

\noindent\textbf{odpowiedzi}:\\ 
\textbf{(a)} $3\pi$, \textbf{(b)} $\frac{31\sqrt{2}}{60}$,  
\textbf{(c)} $\frac{1-\ln 2}{16}$,  \textbf{(d)} $\frac 23$\\

\begin{figure}[H]
	\centering
	\begin{tikzpicture}
		\foreach \x in {0,...,34}
		\draw (0.5*\x,0)--(0.5*\x,19cm);
		\foreach \y in {0,...,38}
		\draw (0,0.5*\y)--(17cm,0.5*\y);
	\end{tikzpicture}
\end{figure}

\begin{figure}[H]
	\centering
	\begin{tikzpicture}
		\foreach \x in {0,...,34}
		\draw (0.5*\x,0)--(0.5*\x,27cm);
		\foreach \y in {0,...,54}
		\draw (0,0.5*\y)--(17cm,0.5*\y);
	\end{tikzpicture}
\end{figure}

\noindent
\textbf{Zadanie 2} \\
\textbf{(a)} wyznacz obj�to�� obszaru ograniczonego powierzchniami $x = 0$, $y = 0$, $z = 0$, $3x + 4y + 6z - 12 = 0$,\\
\textbf{(b)} wyznacz obj�to�� obszaru ograniczonego powierzchniami $z = 0$, $z = 16 - x^2 - y^2$,\\
\textbf{(c)} wyznacz po�o�enie �rodka ci�ko�ci obszaru ograniczonego liniami $y = x$ oraz $y = x^2$,\\ 
\textbf{(d)} wyznacz po�o�enie �rodka ci�ko�ci obszaru $D$: $1\leq x^2 + y^2 \leq 4$, $y \geq x$, $x\geq 0$.\\

\vspace{0.25cm}

\noindent
\textbf{odpowiedzi}:\\ 
\textbf{(a)} $V = 4$, 
\textbf{(b)} $V = 128\pi$, 
\textbf{(c)} $G = \left(\frac 1 2, \frac 2 5\right)$, 
\textbf{(d)} $G\left(\frac{28(2-\sqrt{2})}{9\pi}, \frac{28\sqrt{2}}{9\pi}\right)$.
\begin{figure}[H]
	\centering
	\begin{tikzpicture}
		\foreach \x in {0,...,34}
		\draw (0.5*\x,0)--(0.5*\x,22cm);
		\foreach \y in {0,...,44}
		\draw (0,0.5*\y)--(17cm,0.5*\y);
	\end{tikzpicture}
\end{figure}

\begin{figure}[H]
	\centering
	\begin{tikzpicture}
		\foreach \x in {0,...,34}
		\draw (0.5*\x,0)--(0.5*\x,27cm);
		\foreach \y in {0,...,54}
		\draw (0,0.5*\y)--(17cm,0.5*\y);
	\end{tikzpicture}
\end{figure}

\noindent
\textbf{Zadanie 3} Rozwi�� r�wnania r�niczkowe:\\ 
\textbf{(a)} $\frac{y'}{x^2} - y^2 = 0,\ y(0) = 1$, 
\textbf{(b)} $ y' x^3 - 2y = 0,\ y(1) = \frac 1e$, 
\textbf{(c)} $\frac{y'}{\cos x} + 4\sin x \cos^2 y = 0,\ y\left(\frac{\pi}{4}\right) = 0$, \\
\textbf{(d)} $\frac{y'}{x} - \frac{1}{y^2\sqrt{x^2 + 1}} = 0,\ y(0) = 1$.

\vspace{0.25cm}

\noindent
\textbf{odpowiedzi}:\\ 
\textbf{(a)} $y(x) = \frac{3}{3-x^3}$, 
\textbf{(b)} $y(x) = e^{-\frac{1}{x^2}}$, 
\textbf{(c)} $y(x) = \text{arctg}(\cos(2x))$, 
\textbf{(d)} $y(x) = \sqrt[3]{3\sqrt{x^2+1} - 2}$.

\begin{figure}[H]
	\centering
	\begin{tikzpicture}
		\foreach \x in {0,...,34}
		\draw (0.5*\x,0)--(0.5*\x,23cm);
		\foreach \y in {0,...,46}
		\draw (0,0.5*\y)--(17cm,0.5*\y);
	\end{tikzpicture}
\end{figure}

\begin{figure}[H]
	\centering
	\begin{tikzpicture}
		\foreach \x in {0,...,34}
		\draw (0.5*\x,0)--(0.5*\x,27cm);
		\foreach \y in {0,...,54}
		\draw (0,0.5*\y)--(17cm,0.5*\y);
	\end{tikzpicture}
\end{figure}

\noindent
\textbf{Zadanie 4} Rozwi�� r�wnania r�niczkowe:\\
\textbf{(a)} $3y''-4y'+y = 16 e^{3x}$, 
\textbf{(b)} $y'' + 4y' + 13y = \sin(3x)$, 
\textbf{(c)} $y''-4y'+4y = 2 e^{2x}$, 
\textbf{(d)} $y'' + 4y' + 5y = 25x^2$.

\vspace{0.25cm}

\noindent
\textbf{odpowiedzi}:\\ 
\textbf{(a)} $y(x) = C_1 e^{x/3} + C_2 e^x + e^{3x}$, 
\textbf{(b)} $y(x) = e^{-2x}(C_1\cos(3x) + C_2\sin(3x)) + \frac{1}{40}\sin(3x) - \frac{3}{40}\cos(3x)$,\\ 
\textbf{(c)} $y(x) = C_1 e^{2x} + C_2 x e^{2x} + x^2 e^{2x}$, 
\textbf{(d)} $y(x) = e^{-2x}(C_1\cos x + C_2\sin x) + 5x^2 - 8x + \frac{22}{5}$.

\begin{figure}[H]
	\centering
	\begin{tikzpicture}
		\foreach \x in {0,...,34}
		\draw (0.5*\x,0)--(0.5*\x,23cm);
		\foreach \y in {0,...,46}
		\draw (0,0.5*\y)--(17cm,0.5*\y);
	\end{tikzpicture}
\end{figure}

\begin{figure}[H]
	\centering
	\begin{tikzpicture}
		\foreach \x in {0,...,34}
		\draw (0.5*\x,0)--(0.5*\x,27cm);
		\foreach \y in {0,...,54}
		\draw (0,0.5*\y)--(17cm,0.5*\y);
	\end{tikzpicture}
\end{figure}

\end{document}